\chapter{Unique Positive Scale-Defect Crossing}
\label{chap:unique-scale-defect-crossing}

\status{Proved on the positive real \(u\)-axis. No complex zero-location or RH claim is made.}

For \(u>0\), define
\[
X(u)=\xi\!\left(\frac{u}{1+u}\right),\qquad
\Delta_a(u)=X(au)-X(u),\quad a>1.
\]
Write \(u=e^\lambda\). Then
\[
\frac{u}{1+u}=\frac12+\frac12\tanh\frac\lambda2.
\]
Since \(\xi(1/2+z)=F(z^2)\),
\[
\boxed{X(e^\lambda)=F\!\left(\frac14\tanh^2\frac\lambda2\right).}
\]
From the positive coefficient theorem of Chapter 27,
\[
F(w)=\sum_{k\ge0}a_k w^k,\qquad a_k>0,
\]
so
\[
F'(w)>0\qquad(w\ge0).
\]
Consequently \(X(e^\lambda)\) is even in \(\lambda\) and strictly increasing as a function of \(|\lambda|\).

Let \(L=\log a>0\). Then
\[
\Delta_a(e^\lambda)=X(e^{\lambda+L})-X(e^\lambda),
\]
and therefore its sign is the sign of
\[
|\lambda+L|-|\lambda|.
\]
This gives the following theorem.

\begin{theorem}[Unique positive scale-defect crossing]
For every \(a>1\),
\[
\Delta_a(u)<0\quad\text{for }0<u<a^{-1/2},
\]
\[
\boxed{\Delta_a(u)=0\iff u=a^{-1/2}},
\]
and
\[
\Delta_a(u)>0\quad\text{for }u>a^{-1/2}.
\]
\end{theorem}

For the Uroboros scale \(a=32\), the unique crossing is
\[
\boxed{u=32^{-1/2}}.
\]
This strengthens the antisymmetry result of Chapter 33 on the positive axis: the symmetry-forced crossing is not merely present but unique and sign-changing.

\section{Proof firewall}
The theorem is one-dimensional on the positive real \(u\)-axis. It does not imply PF\(_\infty\), SOH-G003 real-rootedness, complex zero location, or the Riemann Hypothesis.
